\documentclass{axiobac}

\annee{1978}
\mois{Juin}
\session{principale}
\epreuve{Mathématiques (épreuve obligatoire)}
\sectionbac{Math-Sciences et Math-Technique}
\sections{mathematiques, sciences-techniques}
\ministere{Ministère de l'Éducation Nationale}

\begin{document}
\entete

%----------------------------------------------------------------------------------------
%   EXERCICE 1
%----------------------------------------------------------------------------------------
\begin{exercice}{}
\titre{Équation du second degré à paramètre et lieu de points}
\chapitres{Nombres Complexes}
\banque{0000164}

Soit, dans $\mathbb{C}$, l'équation :
\[
z^2 - 2z + 2\sin^2\theta - 2i\sin\theta\cos\theta = 0 \qquad \text{avec } \theta \in [0, \pi[ .
\]

\begin{enumerate}
    \item Résoudre cette équation. \\
    On désignera par $z'$ la racine dont le coefficient $b$ de la partie imaginaire $(ib)$ est positif et par $z''$ l'autre racine.
    \item On pose $z_1 = z' + i\sin\theta$ et $z_2 = z'' - i\sin\theta$. \\
    Soit $M_1$ et $M_2$ les images respectives de $z_1$ et $z_2$ dans le plan complexe.
    \begin{enumerate}
        \item Montrer que $M_1$ et $M_2$ sont symétriques par rapport au point d'affixe $1$.
        \item Déterminer l'ensemble des points $M_1$ lorsque $\theta$ décrit l'intervalle $[0, \pi[$ ; en déduire l'ensemble des points $M_2$.
    \end{enumerate}
\end{enumerate}

\end{exercice}

%----------------------------------------------------------------------------------------
%   EXERCICE 2
%----------------------------------------------------------------------------------------
\begin{exercice}{}
\titre{Somme et écart de deux jetons : loi conjointe}
\chapitres{Probabilités}
\horsprogramme

On dispose de deux boîtes contenant chacune $3$ jetons indiscernables au toucher et numérotés $0$, $1$, $2$. Une épreuve consiste à tirer, au hasard, un jeton de chaque boîte. \\
Soit $X$ l'application qui, à chaque épreuve, associe la somme des nombres inscrits sur les jetons tirés et soit $Y$ l'application qui, à chaque épreuve, associe la valeur absolue de la différence des nombres inscrits sur les jetons tirés.

\begin{enumerate}
    \item Déterminer la loi conjointe du couple $(X, Y)$ et les lois marginales de $X$ et de $Y$.
    \item Calculer les espérances mathématiques de $X$, $Y$ et $X \cdot Y$.
\end{enumerate}

\end{exercice}

%----------------------------------------------------------------------------------------
%   PROBLÈME
%----------------------------------------------------------------------------------------
\begin{probleme}{}

Pour tout couple de réels $(a, b)$, on considère la fonction $f_{(a,b)}$ de $\mathbb{R}^+$ dans $\mathbb{R}$ définie par :
\[
\begin{cases}
    f_{(a,b)}(x) = x^2\,(a + b \operatorname{Log} x) & \text{si } x > 0 \\
    f_{(a,b)}(0) = 0
\end{cases}
\]

\partie{A}
\titre{Étude de la fonction $x \mapsto x^2(1 - \ln x)$}
\chapitres{Fonction logarithme népérien, Dérivabilité, Etude de Fonctions}
\banque{0000165}

On pose $g = f_{(1,-1)}$.

\begin{enumerate}
    \item Montrer que la fonction $g$ est continue au point $0$ et qu'elle est dérivable en ce point.
    \item Étudier les variations de la fonction $g$ et tracer sa courbe représentative dans un plan rapporté à un repère orthonormé. (On donne : $e = 2{,}71\ldots$ ; $\sqrt{e} = 1{,}68\ldots$).
\end{enumerate}

\partie{B}
\titre{Un endomorphisme d'un espace de fonctions}
\horsprogramme

On désigne par $D$ l'espace vectoriel réel des fonctions dérivables sur $\mathbb{R}^+$ et par $E$ l'ensemble des fonctions $f_{(a,b)}$ lorsque $(a, b)$ décrit $\mathbb{R}^2$. On pose :
\[
u = f_{(1,0)} \qquad \text{et} \qquad v = f_{(0,1)} .
\]

\begin{enumerate}
    \item
    \begin{enumerate}
        \item Vérifier que $u$ et $v$ appartiennent à $D$.
        \item Montrer que $E$ est le sous-espace vectoriel de $D$ engendré par $u$ et $v$.
        \item Montrer que $(u, v)$ est une base de $E$.
    \end{enumerate}
    \item On désigne par $\overline{f}_{(a,b)}$ la fonction de $\mathbb{R}^+$ dans $\mathbb{R}$ définie par $\overline{f}_{(a,b)}(x) = x\,f'_{(a,b)}(x)$.
    \begin{enumerate}
        \item Montrer que $\overline{f}_{(a,b)}$ est un élément de $E$.
        \item Soit l'application $\sigma : E \longrightarrow E$, $f_{(a,b)} \longmapsto \overline{f}_{(a,b)}$. \\
        Montrer que $\sigma$ est un endomorphisme de $E$ ; déterminer sa matrice $A$ dans la base $(u, v)$.
    \end{enumerate}
    \item
    \begin{enumerate}
        \item Vérifier que $A^2 - 4A + 4I = O$ ($I$ est la matrice unité, $O$ est la matrice nulle).
        \item En déduire que :
        \[
        \forall x \in \mathbb{R}^*_+, \qquad x^2 f''_{(a,b)}(x) - 3x\,f'_{(a,b)}(x) + 4 f_{(a,b)}(x) = 0 .
        \]
    \end{enumerate}
\end{enumerate}

\partie{C}
\titre{Intégrales de $t^4 \ln t$ et $t^4 (\ln t)^2$, limites en $0$}
\chapitres{Intégrales, Fonction logarithme népérien}
\banque{0000165}

Soit $s$ un réel de l'intervalle $]0, 1[$.

\begin{enumerate}
    \item
    \begin{enumerate}
        \item Calculer $I(s) = \displaystyle\int_s^1 t^4 \operatorname{Log} t \, dt$ ; en déduire $\displaystyle\lim_{s \to 0^+} I(s)$.
        \item On pose $J(s) = \displaystyle\int_s^1 t^4 (\operatorname{Log} t)^2 \, dt$. \\
        Trouver une relation entre $I(s)$ et $J(s)$ ; en déduire $\displaystyle\lim_{s \to 0^+} J(s)$.
    \end{enumerate}
    \item On pose $K(s) = \displaystyle\int_s^1 \left[ f_{(a,b)}(t) \right]^2 dt$. Déterminer $\displaystyle\lim_{s \to 0^+} K(s)$.
\end{enumerate}

\end{probleme}

\end{document}
